\section{Corpus and Orientation Characterization}
\label{sec:corpus}
The comparison asks how documentation serving different purposes organizes
explicit cross-references, and which additional distinctions appear when link
direction is retained.

\begin{table}[t]
\caption{Corpus profiles by documentation purpose. Read LWCC/full-inventory coverage ($n/N$) before comparing magnetic and undirected modes. All eigenvalues use the same LWCC and $W$. Purpose groups overlap; they are descriptive, not measured classes. Similar eigenvalues need not imply similar organization. Numerical zeros are printed as zero.}
\label{tab:corpora}
\centering
\begin{tabular}{@{}lrrrr@{}}
\toprule
& & & \multicolumn{2}{c}{Second-mode comparison}\\
\cmidrule(l){4-5}
Corpus & $n/N$ & $\lambda_1(.25)$ & Magnetic & Undirected\\
& & & $\lambda_2(.25)$ & $\lambda_2(0)$\\
\midrule
\multicolumn{5}{@{}l@{}}{\textit{Operational setup and configuration}}\\
Docker & 844/1263 & 0.0842 & 0.0969 & 0.0101 \\
Kubernetes & 1146/1572 & 0.0504 & 0.1257 & 0.0394 \\
Ansible & 391/492 & 0.0132 & 0.0150 & 0.0037 \\
Argo CD & 393/470 & 0.0038 & 0.0051 & 0.0050 \\
\multicolumn{5}{@{}l@{}}{\textit{Framework learning and reference}}\\
Django & 649/651 & 0.0533 & 0.1873 & 0.0638 \\
Express & 44/52 & 0.0591 & 0.0924 & 0.0854 \\
React & 206/218 & 0.2099 & 0.2237 & 0.1332 \\
Next.js & 295/418 & 0.1296 & 0.1311 & 0.1134 \\
\multicolumn{5}{@{}l@{}}{\textit{Data management and administration}}\\
Elasticsearch & 2257/2364 & 0.0966 & 0.1188 & 0.0271 \\
Redis & 45/97 & 0.0124 & 0.0746 & 0.0161 \\
\multicolumn{5}{@{}l@{}}{\textit{Product workflows and development}}\\
GitHub & 2354/3515 & 0.0139 & 0.0396 & 0.0030 \\
VS Code & 658/770 & 0.0620 & 0.1164 & 0.0611 \\
\multicolumn{5}{@{}l@{}}{\textit{Library APIs and usage}}\\
PyTorch & 9/244 & 0.0000 & 0.2929 & 0.2929 \\
\multicolumn{5}{@{}l@{}}{\textit{Hardware assembly and software setup}}\\
AgOpenGPS & 23/59 & 0.0762 & 0.1527 & 0.1474 \\
\bottomrule
\end{tabular}
\end{table}

\subsection{Documentation Purposes and Connected Coverage}
The panel comprises fourteen frozen retrieval-evaluation collections with
available questions and recoverable supporting documentation.
Table~\ref{tab:corpora} groups their overlapping purposes descriptively;
these are neither measured style classes nor a sample stratified by authoring style.

Connected coverage determines how much of each collection a spectral profile
describes (Table~\ref{tab:corpora}). Django and React have nearly all retained
pages in one component, whereas Docker and GitHub leave substantial portions
outside it. PyTorch's measured component contains only nine of 244 pages and
is a tree, giving a zero lowest eigenvalue at every tested charge. This zero
describes that small extracted component, not the library's documentation as
a whole. Coverage thus distinguishes a broadly connected body from a limited
linked fragment before their eigenvalues are compared. It concerns resolved
in-text links; rendered navigation can connect pages outside this graph.

\subsection{What Retaining Direction Reveals}
Magnetic and undirected measurements describe complementary aspects of each collection.
\begin{table}[t]
\caption{Spearman correlations of magnetic modes with corpus descriptors and the exposure proxy. Orientation variants are dependent; parentheses give within-family Holm-adjusted $p$-values.}
\label{tab:correlations}
\centering
\begin{tabular}{@{}lrr@{}}
\toprule
Comparator & $\lambda_1(.25)$ & $\lambda_2(.25)$\\
\midrule
\multicolumn{3}{@{}l@{}}{\textit{Corpus panel: Spearman $\rho$, 14 anchors}}\\
Matched connectivity $\lambda_2(0)$ & $0.327$ & $0.864$\\
Mean weighted degree & $0.679$ & $0.138$\\
Out-degree Gini & $-0.604$ & $-0.037$\\
Trophic incoherence & $0.442$ & $-0.138$\\
Reciprocity & $0.201$ & $-0.307$\\
\addlinespace
\multicolumn{3}{@{}l@{}}{\textit{Exposure proxy: $\rho$ (Holm $p$), 29 variants/family}}\\
Main & $0.024\;(0.904)$ & $-0.223\;(0.491)$\\
Scaffolding removed & $0.034\;(0.859)$ & $0.244\;(0.394)$\\
\bottomrule
\end{tabular}
\end{table}

Within those components, Docker's second magnetic eigenvalue is substantially
above its matched undirected value, while Argo CD's two values are nearly equal
(Table~\ref{tab:corpora}). Retaining direction therefore changes the low-end
description unevenly across collections. Across charge, React and Next.js rise
comparatively steeply, while Argo CD and Ansible remain low. The matched
reversals in Section~\ref{sec:validation} isolate an orientation response;
these cross-corpus contrasts also involve differences in topology and do not
attribute a response to documentation style.

Similar scores can also conceal different organizations: Ansible and GitHub
have close first eigenvalues despite different sizes and degree concentration.
Across the panel, the second magnetic mode tracks ordinary connectivity more
closely than the first (Table~\ref{tab:correlations}). The first tends to be
higher with greater mean weighted degree and lower with more concentrated
out-degree; reciprocity is more weakly associated. These descriptive pairwise
comparisons support reading the spectrum alongside coverage and simpler
structural measures, rather than assigning a documentation type from one score.

\subsection{Fixed Text, Changed Orientation}
With page content fixed, the orientation analysis compares spectra and evidence access.
With $W$ fixed, the main Docker orientation family has little variation in
$\lambda_1(.25)$; removing scaffolding broadens the response. Full reciprocation makes the ground eigenvalue zero
but changes $W$. Unlike the matched study (Section~\ref{sec:validation}),
the comparisons excluding the fully reciprocal version preserve $W$ but allow directed summaries to change.

The exposure diagnostic expands a fixed semantic seed along outgoing links
into a bounded context, counting added exposure over 142 Docker questions.
This tests whether the orientation differences accompany changes in access to
candidate supporting material.
With automatic, size-biased lexical support candidates, the first mode's
association is close to zero in both families; the second's changes sign
(Table~\ref{tab:correlations}). Neither passes the multiplicity-adjusted screen
for incremental-information analysis.
